% Draft markers for the Method/Experiments drafts (draft-marker-discipline).
% \TODO       an unresolved factual or sourcing question (a number, a run, a protocol)
% \note       editorial remark addressed to the authors
% \directive  a change this passage must receive, and why
% \superseded text replaced by a decision, kept visible in draft mode
% Set \draftmodefalse to hide all four for a clean copy.
% The manuscript preamble already defines \TODO; \providecommand keeps that definition.
\usepackage{xcolor}
\usepackage{multirow}
\newif\ifdraftmode
\draftmodetrue
\providecommand{\TODO}[1]{\textbf{\color{red}[TODO: #1]}}
% \prov{x}   a provisional number: the best value available now, to be replaced by the final runs (source in the
%            % SRC comment next to it). Draft mode prints it in brown; clean mode prints it plainly.
% \pend      a table cell whose run or scoring is still pending: a red dash in draft mode, -- in the clean copy.
\ifdraftmode
  \newcommand{\prov}[1]{{\color{brown}#1}}
  \newcommand{\pend}{{\color{red}--}}
  \newcommand{\note}[1]{{\color{blue}[Note: #1]}}
  \newcommand{\directive}[1]{\par\noindent\fbox{\parbox{0.95\linewidth}{\small\textbf{Directive:} #1}}\par}
  \newcommand{\superseded}[1]{{\color{gray}#1}}
\else
  \newcommand{\prov}[1]{#1}
  \newcommand{\pend}{\nev}
  \renewcommand{\TODO}[1]{}
  \newcommand{\note}[1]{}
  \newcommand{\directive}[1]{}
  \newcommand{\superseded}[1]{}
\fi
